\begin{description}
\item[(6.1)] From Wien's law:
\[ \lambda_{max} = \frac{2.898\times10^{-3}\ \mathrm{m\cdot K}}{T_{eff}\,(\mathrm{K})}
   = \frac{2.898\times10^{-3}\ \mathrm{m\cdot K}}{4995\ \mathrm{K}} \]
\[ \lambda_{max} = 5.8\times10^{-7}\ \mathrm{m} = 580.2\ \mathrm{nm} \]
\hfill(2 points)

\item[(6.2)] With the parallax doable % sic: source wording
to find the distance from Earth:
\[ d = \frac{1}{pllx\,(\mathrm{arcsec})} = \frac{1}{0.035\ \mathrm{arcsec}}
     = 28.55\ \mathrm{pc} \]
\hfill(1 point)
\begin{align*}
M_V &= m_v - 5\log(d\,[\mathrm{pc}]) + 5\\
M_V &= 8.3 - 5\log(28.57) + 5 = 6.02
  % sic: source uses 28.55 pc above but 28.57 here
\end{align*}
\hfill(1 point)

\item[(6.3)] \hfill(2 points)\\
% FIGURE NEEDED: the radial-velocity curve of HD 93083 (HARPS) with two
% horizontal green lines marking the maximum (43.662 km/s) and minimum
% (43.626 km/s) radial velocities (2021_T6_Solution.pdf, page 1).
Each division on Y has 0.002~km\,s$^{-1}$, so the maximum radial velocity is
$43.662$~km\,s$^{-1}$ and the minimus % sic: source misprint for "minimum"
is $43.626$~km\,s$^{-1}$. Then, the mean radial velocity of Macondo is
$v_r = 43.644$~km\,s$^{-1}$.

\item[(6.4)] The tangential velocity of Macondo, from the plot, is its
variation from the mean system velocity
\[ v_s = 43.662\ \mathrm{km\,s^{-1}} - 43.644\ \mathrm{km\,s^{-1}}
       = 0.018\ \mathrm{km\,s^{-1}} \]
The masses of the star and the planet are known so it is only needed to find
orbital velocity of Melquiades. But first it is important to convert the mass
of the star to kg.
\[ 0.7M_\odot \times \frac{1.989\times10^{30}\ \mathrm{kg}}{1M_\odot}
   = 1.4\times10^{30}\ \mathrm{kg} \]
\hfill(1 point)
\[ v_p = \frac{m_s}{m_p}\times v_s
       = \frac{1.4\times10^{30}\ \mathrm{kg}}{7\times10^{26}\ \mathrm{kg}}
         \times 0.02\ \mathrm{km\cdot s^{-1}}
       = 40\ \mathrm{km\cdot s^{-1}} \]
\hfill(1 point)

\item[(6.5)] As the motion of the planet is circular, the orbital period is:
\[ T = \frac{2\pi}{v_p}\times a \]
being $a$ the distance to the central star.

With the 3$^{\mathrm{rd}}$ Kepler's law, the orbital period is
$T^2 = \frac{4\pi^2}{Gm_s}\times a^3$.

Then, combining both expressions of $T$ is possible to find $a$
\begin{align*}
\left(\frac{2\pi}{v_p}\times a\right)^{2} &= \frac{4\pi^2}{Gm_s}\times a^3\\
\left(\frac{2\pi}{v_p}\right)^{2}\frac{Gm_s}{4\pi^2} &= a
\end{align*}
\[ a = 7.25\times10^{10}\ \mathrm{m} = 0.48\ \mathrm{au} \]
\hfill(2 points)

With the previous solution is able to find the orbital period:
\[ T = \frac{2\pi}{v_p}\times a
     = \frac{2\pi\times(7.21\times10^{10}\ \mathrm{m})}{36000\ \mathrm{m\cdot s^{-1}}}
     % sic: source gives a = 7.25e10 m above but uses 7.21e10 m here,
     % and v_p = 40 km/s above but 36000 m/s here
     = 1.27\times10^{7}\ \mathrm{s} = 147\ days \]
\hfill(2 points)
\end{description}
